% Source: tex/chapters/ch05/12_トピックと参考文献の対応.tex
\subsection{トピックと参考文献の対応}
この表は、SWEBOK が示す「トピック対参考資料の行列」を、学習用に簡略化したものである。詳しく学ぶときは、各トピックに対応する文献を読む。

\par\smallskip\noindent\refstepcounter{table}\label{tab:ch05-04}\textbf{表 \thetable: 第05章 ソフトウェアテスト 表04}\par\smallskip
\begin{longtable}{P{0.46\linewidth}P{0.46\linewidth}}
\toprule
\textbf{トピック} & \textbf{主な参考資料} \\
\midrule
1. ソフトウェアテストの基礎 & Naik and Tripathy; Sommerville; Laporte and April \\
1.1 欠陥と障害 & Naik and Tripathy; Sommerville; Laporte and April \\
1.2 主要論点 & Naik and Tripathy; ISO/IEC/IEEE 29119; IEEE 1012; ISO/IEC 25010; ISO/IEC/IEEE 32675 \\
1.3 他活動との関係 & Software Quality KA; Software Construction KA; Software Security KA; Professional Practice KA \\
2. テストレベル & Naik and Tripathy; Sommerville \\
2.1 テスト対象 & Naik and Tripathy; Sommerville \\
2.2 テスト目的 & Naik and Tripathy; Sommerville; ISO/IEC/IEEE 29119; ISO/IEC/IEEE 24765; Nielsen \\
3. テスト技法 & Naik and Tripathy; ISO/IEC/IEEE 29119; Utting et al.; Nielsen \\
3.1 仕様ベース技法 & Naik and Tripathy; ISO/IEC/IEEE 29119; Utting et al. \\
3.2 構造ベース技法 & Naik and Tripathy; ISO/IEC/IEEE 29119 \\
3.3 経験ベース技法 & ISO/IEC/IEEE 29119; Riccio et al. \\
3.4 欠陥ベース・ミューテーション & Naik and Tripathy; Papadakis et al.; ISO/IEC/IEEE 24765 \\
3.5 利用ベース技法 & Naik and Tripathy; Sommerville; Nielsen \\
3.6 アプリケーションの性質 & Sommerville; Laporte and April; IEEE 1012 \\
3.7 技法の選択と組合せ & Laporte and April; ISO/IEC/IEEE 32675; ISO/IEC/IEEE 29119 \\
3.8 派生知識に基づく技法 & Sommerville; Laporte and April; Utting et al. \\
4. テスト関連測度 & Sommerville; Laporte and April; ISO/IEC/IEEE 32675; ISO/IEC/IEEE 29119 \\
5. テストプロセス & Sommerville; ISO/IEC/IEEE 29119; SWECOM; ISO/IEC 20246 \\
6. 開発プロセスと適用領域 & Sommerville; Laporte and April; ISO/IEC/IEEE 29119 \\
7. 新興技術のテストと新興技術によるテスト & Riccio et al.; Demi et al.; Bertolino et al.; Achary and Raj; Mao et al. \\
8. ソフトウェアテストツール & Naik and Tripathy; Laporte and April \\
\bottomrule
\end{longtable}
